% ======================================================================
\section{Setup: soft gluon wave function and the evolution operator}\label{sec:setup}
% ======================================================================

\purpose{The whole calculation is an evaluation of products of the coefficients of one operator, $\Omega$, which dresses the valence charges of the projectile with soft gluons. This section defines $\Omega$, explains how its coefficients are obtained order by order from light-cone perturbation theory and from unitarity, says why one coefficient of order $g^3$ is needed and where its ingredients come from, and lists all coefficients in the form in which they are used. The one conceptual point that goes beyond a standard perturbative calculation is that the couplings to the valence charges are operators, so that the order of the factors has to be kept.}

In this section we summarise the framework of~\cite{Radhakrishnan:2026soft}, on which the present calculation is built, and collect the coefficients of the evolution operator $\Omega$ that enter the single inclusive cross section. The coefficients of order $g$ and $g^2$ were derived in~\cite{Radhakrishnan:2026soft}; the coefficient of order $g^3$ that creates one gluon, $\mf A^{(3)}_2$, is obtained here. The derivations of the light-cone wave functions from which the coefficients are obtained are summarised in Appendix~\ref{app:wavefunctions}.

\subsection{Light-cone conventions}\label{sec:conventions}

We use light-cone coordinates $x^\pm=(x^0\pm x^3)/\sqrt2$ and denote transverse vectors by $\v x=(x^1,x^2)$; where no confusion is possible the transverse vectors are written without boldface, and $\int_{x}\equiv\int d^2x$. The projectile moves in the $+$ direction, $x^+$ plays the role of time, and $p^+>0$ is the longitudinal momentum of a physical particle. The light-cone energy of an on-shell gluon is $p^-=\v p^2/2p^+$, and we write
\begin{align}
&E_g(p)=\frac{\v p^2}{2p^+}\,,\qquad E_{gg}(p,q)=\frac{\v p^2}{2p^+}+\frac{\v q^2}{2q^+}\,,\notag\\ &d^3p\equiv dp^+d^2p\,,\qquad \delta^{(3)}(p-q)\equiv\delta(p^+-q^+)\,\delta^{(2)}(\v p-\v q)\,.
\end{align}
We work in the light-cone gauge $A^+=0$, in which the transverse components $A_i^a$, $i=1,2$, are the dynamical fields. The gluon creation and annihilation operators satisfy
\begin{align}\label{eq:aadagger}
\big[a^a_i(k),a^{\dagger b}_j(p)\big]=(2\pi)^3\,\delta^{ab}\,\delta_{ij}\,\delta^{(3)}(k-p)\,.
\end{align}
For the Fourier transforms in the transverse plane we use the conventions
\begin{align}\label{eq:FTconventions}
&a(k^+,\v k)=\int d^2z\,e^{-i\v k\cdot\v z}\,a(k^+,\v z)\,,\qquad
a^\dagger(k^+,\v k)=\int d^2z\,e^{i\v k\cdot\v z}\,a^\dagger(k^+,\v z)\,,\notag\\
&f(\v k)=\int d^2x\,e^{i\v k\cdot\v x}\,f(\v x)\,,
\end{align}
so that $\big[a^a_i(k^+,\v z),a^{\dagger b}_j(p^+,\v z')\big]=2\pi\,\delta^{ab}\delta_{ij}\,\delta(k^+-p^+)\,\delta^{(2)}(\v z-\v z')$. One- and two-gluon states are normalised as
\begin{align}\label{eq:states}
\ket{g^a_i(k^+,\v z)}=\frac{a^{\dagger a}_i(k^+,\v z)}{(2\pi)^{1/2}}\ket{0}\,,\qquad
\ket{g^a_i(k^+,\v z)\,g^b_j(p^+,\v z')}=\frac{a^{\dagger a}_i(k^+,\v z)\,a^{\dagger b}_j(p^+,\v z')}{2\pi}\ket{0}\,,
\end{align}
with $\braket{g^b_j(p^+,\v z')|g^a_i(k^+,\v z)}=\delta^{ab}\delta_{ij}\,\delta(k^+-p^+)\,\delta^{(2)}(\v z-\v z')$; in momentum space $\ket{g^a_i(k)}=(2\pi)^{-3/2}a^{\dagger a}_i(k)\ket0$ (Appendix~\ref{app:states}).

Throughout the paper we use the Weizs\"acker--Williams field of a point charge and the transverse kernels built from it,
\begin{align}\label{eq:kernels}
\Kv^i(\v x)\equiv\frac{x^i}{\v x^2}\,,\qquad
K(\v x,\v y,\v z)\equiv\Kv(\v x-\v z)\cdot\Kv(\v y-\v z)=\frac{(\v x-\v z)\cdot(\v y-\v z)}{(\v x-\v z)^2(\v y-\v z)^2}\,,
\end{align}
so that $K(\v x,\v x,\v z)=1/(\v x-\v z)^2$ and $K(\v x,\v x,\v z)-2K(\v x,\v y,\v z)+K(\v y,\v y,\v z)=(\v x-\v y)^2/[(\v x-\v z)^2(\v y-\v z)^2]$ is the dipole kernel.

\subsection{Soft and valence modes; the eikonal Hamiltonian}\label{sec:hamiltonian}

Following the CGC logic~\cite{McLerran:1993ni,McLerran:1993ka,Jalilian-Marian:1997jx,Kovner:2005en}, the gluon modes of the projectile are divided according to their longitudinal momentum. The valence modes have $k^+>\vee$, the soft modes have
\begin{align}
\Lambda<k^+<\vee\,,\qquad \vee\equiv\Lambda\,e^{\delta Y}\,,
\end{align}
i.e.\ $\vee$ is the upper and $\Lambda$ the lower cut-off on the plus momentum of the soft modes, and the modes with $k^+<\Lambda$ do not belong to the projectile wave function; their effect is part of the target. The rapidity interval $\delta Y=\log(\vee/\Lambda)$ is the interval over which the soft modes are integrated; the notation is that of~\cite{Lublinsky:2016meo,Radhakrishnan:2026soft}. Since the valence modes are much faster than the soft ones, they act on the soft modes only through their colour charge density
\begin{align}\label{eq:rho}
\rho^a(\v x)\,,\qquad \big[\rho^a(\v x),\rho^b(\v y)\big]=i f^{abc}\,\rho^c(\v x)\,\delta^{(2)}(\v x-\v y)\,,\qquad
\big[\rho^a(\v k),\rho^b(\v p)\big]=i f^{abc}\,\rho^c(\v k+\v p)\,,
\end{align}
which is an operator on the valence Hilbert space; its microscopic definition in terms of the valence gluon operators is given in~\cite{Radhakrishnan:2026soft}. We keep the charges non-commuting throughout. This is necessary: products of $\rho$'s appear in the cross section in a definite order, determined by the time ordering of the light-cone perturbation theory, and the commutators~\eqref{eq:rho} reduce the number of $\rho$'s by one. These reductions produce some of the logarithmic terms of Section~\ref{sec:rapidity}.

After the decomposition of the gauge field into valence and soft components, all terms suppressed by powers of $k^+_{\rm soft}/k^+_{\rm valence}$ are dropped (eikonal approximation). The light-cone Yang--Mills Hamiltonian ($N_f=0$) then takes the form~\cite{Lublinsky:2016meo,Radhakrishnan:2026soft}
\begin{align}\label{eq:Hdecomp}
&H=H_0+H_g+H_{ggg}+H_{gg\text{-inst}}+H_{gggg}+\dots,\notag\\
&H_0=\int_\Lambda^\vee\frac{dk^+}{2\pi}\int\frac{d^2k}{(2\pi)^2}\,\frac{\v k^2}{2k^+}\,a^{\dagger a}_i(k^+,\v k)\,a^a_i(k^+,\v k)\,,
\end{align}
where the dots stand for terms which do not contribute at the order considered. The coupling of the soft gluons to the valence charges is
\begin{align}\label{eq:Hg}
H_g=\int_\Lambda^\vee\frac{dk^+}{2\pi}\int\frac{d^2k}{(2\pi)^2}\,\frac{g\,k^i}{\sqrt2\,|k^+|^{3/2}}\Big[a^{\dagger a}_i(k^+,\v k)\,\rho^a(-\v k)+a^a_i(k^+,\v k)\,\rho^a(\v k)\Big]\,.
\end{align}
It describes the emission and absorption of a soft gluon by the valence charges. $H_{ggg}$ is the triple-gluon vertex among the soft modes, $H_{gg\text{-inst}}$ the instantaneous (Coulomb-like) interaction between two soft gluons and the valence charge density, and $H_{gggg}$ the four-gluon interaction of the soft modes. Their explicit forms and the matrix elements used in this paper are listed in Appendix~\ref{app:hamiltonian}.

\subsection{Light-cone perturbation theory with operator-valued matrix elements}\label{sec:lcpt}

The soft gluon wave function of the projectile is obtained in time-independent (old-fashioned) perturbation theory in $H-H_0$ on the light cone~\cite{Lepage:1980fj,Brodsky:1997de}. Since the matrix elements of $H_g$ and $H_{gg\text{-inst}}$ are operators on the valence space, the perturbative series is written with the matrix elements in the order of the light-cone time ordering: the vertex acting later stands to the left. With $E_0=0$ for the soft vacuum, the first terms of the expansion of the dressed state in terms of free Fock states are
\begin{align}\label{eq:lcpt}
\ket{\Psi}=\ &\ket0-\sum_n\ket n\frac{V_{n0}}{E_n}+\sum_{n,m}\ket n\frac{V_{nm}V_{m0}}{E_nE_m}
-\sum_{n,m,l}\ket n\frac{V_{nm}V_{ml}V_{l0}}{E_nE_mE_l}
+\sum_{n,m}\ket n\frac{V_{n0}\,V_{0m}V_{m0}}{E_n^2E_m}
\notag\\
&-\frac12\sum_m\frac{V_{0m}V_{m0}}{E_m^2}\ket0
+\frac12\sum_{n,m}\ket n\frac{V_{n0}\,V_{0m}V_{m0}}{E_nE_m^2}+\dots,
\end{align}
where $V_{nm}=\braket{n|H-H_0|m}$ and the sums run over Fock states of the soft sector. The first line is the Rayleigh--Schr\"odinger series with intermediate normalisation; the second line implements the canonical normalisation of the state and is needed at orders $g^2$ (vacuum term) and $g^3$ (one-gluon term). 

\paragraph{Nuance: the order of operator-valued matrix elements.} In ordinary perturbation theory the matrix elements $V_{nm}$ are numbers and their order in a product is irrelevant. Here $V_{n0}\propto\rho$ and $V_{0m}V_{m0}\propto\rho\rho$ are operators on the valence space that do not commute with each other, and each term of Eq.~\eqref{eq:lcpt} has a definite order: the factors are written in the order in which the interactions act, the later one to the left. For the Rayleigh--Schr\"odinger terms of the first line this is the familiar time ordering. For the normalisation terms the order is less obvious, and it matters. In the last term of each line the single-emission matrix element $V_{n0}$ stands to the left of the norm factor $V_{0m}V_{m0}$; this ordering follows from the perturbative construction of the unitary operator that diagonalises $H$ with operator-valued couplings, and we have verified it against an explicit construction of $\Omega$ with finite-dimensional matrices representing the charges. Energy denominators that vanish inside an integration region are understood as principal values (Appendix~\ref{app:distributions}); this prescription is used throughout (Section~\ref{sec:fin_conv}).

\subsection{The evolution operator}\label{sec:omega}

The dressed (physical) states of the projectile are obtained from the free Fock states by a unitary operator $\Omega$, $\Omega^\dagger\Omega=1$. Since the soft modes are treated in the eikonal approximation, $\Omega$ can be written as a normal-ordered series in the soft gluon operators, with coefficients that are functionals of $\rho$~\cite{Radhakrishnan:2026soft}:
\begin{align}\label{eq:OmegaExp}
\Omega = \ &\mf{N}
+\int\frac{d^3p}{(2\pi)^3}\Big[\mf{A}^b_{1j}(p)\,a^b_j(p)+\mf{A}^b_{2j}(p)\,a^{\dagger b}_j(p)\Big]
\notag\\&
+\int\frac{d^3p}{(2\pi)^3}\frac{d^3q}{(2\pi)^3}\Big[\mf{B}^{cb}_{1kj}(p,q)\,a^b_j(p)a^c_k(q)+\mf{B}^{cb}_{2kj}(p,q)\,a^{\dagger b}_j(p)a^{\dagger c}_k(q)
\notag\\&\hspace{3.4cm}+\mf{B}^{cb}_{3kj}(p,q)\,a^{\dagger c}_k(q)a^b_j(p)\Big]
\notag\\&
+\int\frac{d^3p}{(2\pi)^3}\frac{d^3q}{(2\pi)^3}\Big[\mf{C}^{dcb}_{1lkj}(p,q)\,a^{\dagger d}_l(q)a^{\dagger c}_k(p-q)a^b_j(p)
\notag\\&\hspace{3.4cm}+\mf{C}^{dcb}_{2lkj}(p,q)\,a^{\dagger b}_j(p)a^c_k(p-q)a^d_l(q)\Big]+\dots
\end{align}
The dots denote terms with more gluon operators, which do not contribute to the single inclusive cross section at $O(g^4)$. The coefficients have definite orders in $g$: $\mf N=1+\mf N_2$ with $\mf N_2=O(g^2)$; $\mf A_{1,2}=\mf A^{(1)}_{1,2}+\mf A^{(3)}_{1,2}$ with $\mf A^{(n)}=O(g^n)$; $\mf B_{1,2,3}=O(g^2)$; $\mf C_{1,2}=O(g)$. The coefficients that multiply only creation operators, $\mf A_2$ and $\mf B_2$, together with $\mf N$, are fixed by the action of $\Omega$ on the soft vacuum, i.e.\ by the light-cone wave function of the projectile; $\mf B_3$ and $\mf C_1$ are fixed by the action of $\Omega$ on a one-gluon state. The remaining coefficients follow from unitarity, imposed order by order in $g$ (Appendix~\ref{app:unitarity}):
\begin{align}\label{eq:unitarity_short}
&\mf A^{(1)\dagger b}_{1j}(p)=-\mf A^{(1)b}_{2j}(p)\,,\qquad
\mf C^{\dagger dcb}_{2lkj}(p,q)=-\mf C^{dcb}_{1lkj}(p,q)\,,\notag\\
&\mf N=1-\frac12\int\frac{d^3p}{(2\pi)^3}\,\mf A^{(1)b}_{1j}(p)\,\mf A^{(1)\dagger b}_{1j}(p)\,,
\end{align}
together with the relations for $\mf A^{(3)}_1$ and $\mf B_1$ given in Eqs.~\eqref{eq:A31unit} and~\eqref{eq:B1unit}. In the relation for $\mf A^{(3)}_1$ the operator character of the coefficients matters: $\mf N_2$ does not commute with $\mf A^{(1)}$, since $[\rho^a(\v x)\rho^a(\v y),\rho^b(\v w)]\neq0$ for $\v x\ne\v y$.

In transverse coordinate space, after the unitarity relations~\eqref{eq:unitarity_short} are used, the part of $\Omega$ that contributes is
\begin{align}\label{eq:OmegaX}
\Omega = \ &\mf N
+\int\frac{dp^+}{2\pi}\int_x\Big[\mf A^{(1)b}_{1j}(p^+,\v x)\,a^b_j(p^+,\v x)-\mf A^{\dagger(1)b}_{1j}(p^+,\v x)\,a^{\dagger b}_j(p^+,\v x)
\notag\\&\hspace{2.6cm}
+\mf A^{(3)b}_{1j}(p^+,\v x)\,a^b_j(p^+,\v x)+\mf A^{(3)b}_{2j}(p^+,-\v x)\,a^{\dagger b}_j(p^+,\v x)\Big]
\notag\\&
+\int\frac{dp^+}{2\pi}\frac{dq^+}{2\pi}\int_{x,y}\Big[\mf B^{cb}_{1kj}(p^+,\v x;q^+,\v y)\,a^b_j(p^+,\v x)a^c_k(q^+,\v y)
\notag\\&\hspace{3.2cm}
+\mf B^{cb}_{2kj}(p^+,-\v x;q^+,-\v y)\,a^{\dagger b}_j(p^+,\v x)a^{\dagger c}_k(q^+,\v y)
\notag\\&\hspace{3.2cm}
+\mf B^{cb}_{3kj}(p^+,\v x;q^+,-\v y)\,a^{\dagger c}_k(q^+,\v y)a^b_j(p^+,\v x)\Big]
\notag\\&
+\int\frac{dp^+}{2\pi}\frac{dq^+}{2\pi}\int_{x,y,z}\Big[\mf C^{dcb}_{1lkj}(p^+,\v z-\v y;q^+,\v y-\v x)
\notag\\&\hspace{3.2cm}\times a^{\dagger d}_l(q^+,\v x)a^{\dagger c}_k(p^+-q^+,\v y)a^b_j(p^+,\v z)
\notag\\&\hspace{3.2cm}
-\mf C^{\dagger dcb}_{1lkj}(p^+,\v z-\v y;q^+,\v y-\v x)
\notag\\&\hspace{3.2cm}\times a^{\dagger b}_j(p^+,\v z)a^c_k(p^+-q^+,\v y)a^d_l(q^+,\v x)\Big]+\dots
\end{align}
The transverse arguments of the coefficients are fixed by the Fourier conventions~\eqref{eq:FTconventions}.

\subsection{How \texorpdfstring{$\Omega$}{Omega} is built order by order}\label{sec:construction}

The coefficients of Eq.~\eqref{eq:OmegaExp} are not postulated; they are determined by two requirements, and it is useful to see how these work before the results are listed.

\paragraph{Requirement 1: $\Omega$ reproduces the perturbative eigenstates.} $\Omega$ maps the free soft vacuum onto the dressed vacuum and a free one-gluon state onto the dressed one-gluon state. When the normal-ordered expansion~\eqref{eq:OmegaExp} acts on $\ket0$, every term that contains an annihilation operator vanishes, and only the terms built from creation operators survive,
\begin{align}\label{eq:OmegaVacShort}
\Omega\ket0=\mf N\ket0+\int\frac{d^3p}{(2\pi)^3}\,\mf A^b_{2j}(p)\,a^{\dagger b}_j(p)\ket0+\int\frac{d^3p}{(2\pi)^3}\frac{d^3q}{(2\pi)^3}\,\mf B^{cb}_{2kj}(p,q)\,a^{\dagger b}_j(p)a^{\dagger c}_k(q)\ket0+\dots
\end{align}
(Eq.~\eqref{eq:Omegavac} with the normalisation of the states written out). Comparing this, Fock sector by Fock sector, with the perturbative state~\eqref{eq:lcpt} fixes $\mf N$, $\mf A_2$ and $\mf B_2$: the zero-gluon sector gives $\mf N$, the one-gluon sector $\mf A_2$, the two-gluon sector $\mf B_2$. When $\Omega$ acts on a one-gluon state $\ket{g^a_i(k)}$, the terms with one annihilation operator contribute as well: $\mf A_1a$ leads back to the vacuum, $\mf B_3a^\dagger a$ to a one-gluon state and $\mf C_1a^\dagger a^\dagger a$ to a two-gluon state, Eq.~\eqref{eq:Omegag}. Comparing with the perturbative dressing of the one-gluon state, in which the energy denominators are $E_g(k)-E_n$ instead of $-E_n$, fixes $\mf B_3$ and $\mf C_1$.

\paragraph{Requirement 2: unitarity.} The coefficients that multiply annihilation operators only, $\mf A_1$ and $\mf B_1$, and $\mf C_2$, are not fixed by the action on the vacuum or on a one-gluon state at the order needed. They follow from $\Omega^\dagger\Omega=1$. The lowest order shows the mechanism. With $\Omega=1+\int\frac{d^3p}{(2\pi)^3}[\mf A^{(1)}_1a+\mf A^{(1)}_2a^\dagger]+O(g^2)$,
\begin{align}
\Omega^\dagger\Omega=1+\int\frac{d^3p}{(2\pi)^3}\Big[\big(\mf A^{(1)b}_{1j}+\mf A^{(1)\dagger b}_{2j}\big)a^b_j+\big(\mf A^{(1)b}_{2j}+\mf A^{(1)\dagger b}_{1j}\big)a^{\dagger b}_j\Big]+O(g^2)\,,
\end{align}
so that $\mf A^{(1)\dagger}_1=-\mf A^{(1)}_2$. At order $g^2$ the coefficient of the identity receives $\mf N_2+\mf N_2^\dagger$ and, from the normal ordering of $\mf A^{(1)\dagger}_2a\times\mf A^{(1)}_2a^\dagger$ with the commutator~\eqref{eq:aadagger}, the contraction $\int\frac{d^3p}{(2\pi)^3}\mf A^{(1)\dagger b}_{2j}(p)\mf A^{(1)b}_{2j}(p)$; hence
\begin{align}\label{eq:N2fromunit}
\mf N_2=-\frac12\int\frac{d^3p}{(2\pi)^3}\,\mf A^{(1)\dagger b}_{2j}(p)\,\mf A^{(1)b}_{2j}(p)
=-\frac{g^2}{8\pi^3}\int_\Lambda^\vee\frac{dp^+}{p^+}\int d^2p\,\frac{\rho^b(\v p)\rho^b(-\v p)}{\v p^2}\,,
\end{align}
where $\mf A^{(1)b}_{2j}(p)=-\sqrt2\,g\,p^j\rho^b(-\v p)/(\sqrt{p^+}\,\v p^2)$ (Appendix~\ref{app:wavefunctions}) was used. This is the first relation of Eq.~\eqref{eq:unitarity_short} and reproduces Eq.~\eqref{eq:N2} below: the normalisation of the state is not an independent input but is forced by unitarity, and it is minus one half of the total probability of emitting one soft gluon. In the same way the coefficients of $a$, $aa$ and $a^\dagger aa$ at higher orders give $\mf A^{(3)}_1$, $\mf B_1$ and $\mf C_2$ (Appendix~\ref{app:unitarity}), and the coefficient of $a^\dagger a$ gives a relation that is not needed for the cross section and serves as a check of $\mf B_3$.

\paragraph{Order by order.} Table~\ref{tab:coeffs} summarises which coefficient appears at which order, what fixes it, which diagrams of Appendix~\ref{app:wavefunctions} contribute to it, and in which rows of the cross section it is used.
\begin{table}[t]
\centering
\renewcommand{\arraystretch}{1.25}
\footnotesize\setlength{\tabcolsep}{4pt}
\begin{tabular}{@{}l c >{\raggedright\arraybackslash}p{0.29\textwidth} >{\raggedright\arraybackslash}p{0.16\textwidth} >{\raggedright\arraybackslash}p{0.21\textwidth}@{}}
\toprule
coefficient & order & fixed by & diagrams & enters rows\\
\midrule
$\mf A^{(1)}_2\equiv\mf A$ & $g$ & $\Omega\ket0$, one gluon; first-order LCPT with $H_g$ & Figure~\ref{fig:vac_g} & LO and all rows\\
$\mf A^{(1)}_1$ & $g$ & unitarity, $=-\mf A^{(1)\dagger}_2$ & --- & eliminated\\
$\mf C_1$ & $g$ & $\Omega\ket g$, two gluons; first-order LCPT with $H_{ggg}$ & Figure~\ref{fig:g_gg} & I.V, II.IV, III.III, III.V, III.VI\\
$\mf C_2$ & $g$ & unitarity, $=-\mf C_1^\dagger$ & --- & eliminated\\
$\mf N_2$ & $g^2$ & norm of $\Omega\ket0$, equivalently unitarity & --- & I.I, I.II (through $\mf A\mf N_2$)\\
$\mf B_2$ & $g^2$ & $\Omega\ket0$, two gluons; second-order LCPT & Figure~\ref{fig:vac_gg} & I.IV, I.V, II.I, III.I--III.III\\
$\mf B_3$ & $g^2$ & $\Omega\ket g$, one gluon; second-order LCPT, denominators $E_g(k)-E_n$ & Figure~\ref{fig:g_g} & I.III\\
$\mf B_1$ & $g^2$ & unitarity, Eq.~\eqref{eq:B1unit} & --- & eliminated\\
$\mf A^{(3)}_2\equiv\mf A^{(3)}$ & $g^3$ & $\Omega\ket0$, one gluon; third-order LCPT & Figs.~\ref{fig:vac_g3_onerho}--\ref{fig:vac_g3_threerho} & I.II\\
$\mf A^{(3)}_1$ & $g^3$ & unitarity, Eq.~\eqref{eq:A31unit} & --- & eliminated; leaves $-[\bar{\mf N}_2,\bar{\mf A}]$\\
\bottomrule
\end{tabular}
\caption{The coefficients of $\Omega$ needed for the single inclusive spectrum at $O(g^4)$. ``Eliminated'' means that the coefficient is expressed through the others by a unitarity relation before the cross section is evaluated.}
\label{tab:coeffs}
\end{table}

\subsubsection*{Why \texorpdfstring{$\mf A^{(3)}_2$}{A2(3)} is needed, and where its ingredients come from}\label{sec:whyA3}

The LO amplitude is of order $g$, so that the LO spectrum is of order $g^2$. At order $g^4$ the spectrum contains (i) the squares and interferences of the $O(g^2)$ amplitudes with two gluons in the final state, which need only the coefficients of order $g$ and $g^2$ of~\cite{Radhakrishnan:2026soft}, and (ii) the interference of the $O(g)$ amplitude with the $O(g^3)$ amplitude with \emph{one} gluon in the final state. The latter contains every way in which $\bar\Omega^\dagger\Omega$ connects the soft vacuum with a one-gluon state at order $g^3$: products of lower-order coefficients ($\mf B_3^\dagger\mf A$, $\mf A^\dagger\mf B_2$, $\mf C_1^\dagger\mf B_2$, $\mf A\mf N_2$), and the genuinely new term in which $\Omega$ itself creates one gluon from the vacuum at order $g^3$. This is $\mf A^{(3)}_2$; it is the only coefficient of the calculation that contains a transverse loop integral with an ultraviolet divergence, the self-energy of the emitted gluon.

By Eq.~\eqref{eq:lcpt} the one-gluon component of the dressed vacuum at order $g^3$ consists of three terms,
\begin{align}\label{eq:thirdorder}
-\frac{V_{nm}V_{ml}V_{l0}}{E_nE_mE_l}\,,\qquad+\frac{V_{n0}\,V_{0m}V_{m0}}{E_n^2E_m}\,,\qquad+\frac12\,\frac{V_{n0}\,V_{0m}V_{m0}}{E_nE_m^2}\,:
\end{align}
the third-order Rayleigh--Schr\"odinger term and the two terms that keep the state normalised. The interaction vertices that can occur are $H_g$ (one charge), $H_{ggg}$ (no charge) and $H_{gg\text{-inst}}$ (one charge, order $g^2$). Sorting the diagrams by the number of valence vertices gives the three parts of $\mf A^{(3)}_2$:
\begin{itemize}
\item \emph{One valence vertex} ($H_g$, then $H_{ggg}$ twice): the emitted gluon fluctuates into a pair and recombines (Figure~\ref{fig:vac_g3_onerho}). This is the self-energy of the emitted gluon. It is linear in $\rho$ and contains the ultraviolet pole, the coefficient $\frac{11}{3}$ and the constant $\frac{67}{9}-\frac{\pi^2}{3}$; the origin of these numbers in the three longitudinal integrals~\eqref{eq:xiintegrals} is shown in Section~\ref{sec:ex_I_II}.
\item \emph{Two valence vertices} ($H_g$ twice with one $H_{ggg}$, or $H_g$ with $H_{gg\text{-inst}}$; Figure~\ref{fig:vac_g3_tworho}). The product of the two charges is split into its commutator and anticommutator, Eq.~\eqref{eq:decomp}. The commutator is one charge and adds to the one-$\rho$ part; the anticommutator gives the two-$\rho$ part $\propto f^{abc}\{\rho^b,\rho^c\}$, Eq.~\eqref{eq:A3tworho}.
\item \emph{Three valence vertices} ($H_g$ three times; Figure~\ref{fig:vac_g3_threerho}), together with the two normalisation-type terms of Eq.~\eqref{eq:lcpt}: a second soft gluon is emitted and reabsorbed by the valence charges. Their commutators feed the one- and two-$\rho$ parts, and what remains is the three-$\rho$ part $\mf A^{(1)}_2\mf N_2$, Eq.~\eqref{eq:A3threerho}.
\end{itemize}
Thus the same diagram contributes to parts with different numbers of charges: the classification by the number of $\rho$'s is made \emph{after} the products of charges have been decomposed, and it is this decomposition that makes the one-$\rho$ part depend on $\vee$ as well as on $\Lambda$ (the terms $\log(1-k^+/\vee)$ of Eq.~\eqref{eq:A3onerho} come from the commutator parts of the diagrams with two and three valence vertices, in which a parent gluon of momentum $k^+/\xi<\vee$ transfers longitudinal momentum to the valence charges; Appendix~\ref{app:wavefunctions}).

\paragraph{Interpretation.} $\Omega\ket0$ is the cloud of soft gluons that accompanies the valence charges. To lowest order it is the Weizs\"acker--Williams field, $\mf A$; $\mf N_2$ is the corresponding loss of norm of the bare component; $\mf B_2$ contains two independent quanta of this field and, in addition, the first effect of the gluon self-interaction, a quantum that has split into two; $\mf A^{(3)}$ is the correction to the one-quantum component due to a second, virtual, soft gluon. $\mf B_3$ and $\mf C_1$ do not describe the cloud of the valence charges but what happens to a soft gluon that is already present, and they enter the observable because the measured gluon is a dressed gluon (Section~\ref{sec:observable_def}).

\subsection{The coefficients}\label{sec:coefficients}

We now list the coefficients in transverse coordinate space. They are obtained by comparing $\Omega\ket0$ and $\Omega\ket{g^a_i(k^+,\v z)}$ with the LCPT wave functions of Appendix~\ref{app:wavefunctions}, sector by sector in Fock space; the factors of $2\pi$ come from the normalisation~\eqref{eq:states}. The momentum-space expressions and the intermediate steps are given in Appendix~\ref{app:wavefunctions} (for the coefficients of order $g$ and $g^2$ also in~\cite{Radhakrishnan:2026soft}).

\paragraph{Normalisation.} The coefficient of the identity is fixed by the norm of the LO wave function,
\begin{align}\label{eq:N2}
\mf N=1+\mf N_2\,,\qquad
\mf N_2&=-\frac{g^2}{8\pi^3}\,\log\frac{\vee}{\Lambda}\int_{x,y,z}K(\v x,\v y,\v z)\,\rho^a(\v x)\rho^a(\v y)
\notag\\&=-\frac{g^2}{8\pi^3}\,\log\frac{\vee}{\Lambda}\int d^2p\,\frac{\rho^a(-\v p)\rho^a(\v p)}{\v p^2}\,.
\end{align}
$\mf N_2$ is the (virtual) probability of not emitting a soft gluon anywhere in the interval $\Lambda<p^+<\vee$, and it carries the full logarithm $\log(\vee/\Lambda)$.

\paragraph{Single-gluon emission at $O(g)$.} The Weizs\"acker--Williams amplitude for the emission of a soft gluon at $\v z$ is
\begin{align}\label{eq:A1}
\mf A^{(1)a}_{2i}(k^+,-\v z)=-\mf A^{(1)\dagger a}_{1i}(k^+,\v z)=\frac{ig}{\sqrt2\,\pi\sqrt{k^+}}\int_x\Kv^i(\v x-\v z)\,\rho^a(\v x)\,.
\end{align}
It corresponds to the LO wave function $\ket{\Psi^{\rm LO}}=\frac{ig}{2\pi^{3/2}}\int_\Lambda^\vee\frac{dk^+}{\sqrt{k^+}}\int_{x,z}\Kv^i(\v x-\v z)\,\rho^a(\v x)\ket{g^a_i(k^+,\v z)}$, Eq.~\eqref{eq:PsiLO}. In what follows we use the shorthand $\mf A\equiv\mf A^{(1)}_2$ when no confusion can arise, and a bar denotes the rotation of all charges by the target, $\bar{\mf A}^a(\v x)=\mf A^a(\v x)\big|_{\rho(\v u)\to U(\v u)\rho(\v u)}$, see Section~\ref{sec:observable}.

\paragraph{Single-gluon emission at $O(g^3)$.} The coefficient $\mf A^{(3)}_2$ contains terms with one, two and three charges. In momentum space the one-$\rho$ part reads, for $\Lambda\to0$ wherever this is possible,
\begin{align}\label{eq:A3onerho}
\Big[\mf A^{(3)a}_{2i}(k^+,\v k)\Big]_{\rho}=\ &\frac{g^3N_c\,\rho^a(-\v k)\,k^i}{8\sqrt{2k^+}\,\pi^2\,\v k^2}\Bigg\{\Big[-\frac2\epsilon+\log\frac{\v k^2}{\mu^2_{\MS}}\Big]
\notag\\&\hspace{1.5cm}\times\Big[\frac{11}{3}-2\log\frac{k^+}{\Lambda}+\log\Big(1-\frac{k^+}{\vee}\Big)\Big]
\notag\\&\hspace{1.5cm}
+\frac{\pi^2}{3}-\frac{67}{9}+\log^2\Big(1-\frac{k^+}{\vee}\Big)+\log\Big(1-\frac{k^+}{\vee}\Big)\log\frac{\vee}{k^+}\Bigg\},
\end{align}
where the transverse integrals are performed in $d=2-\epsilon$ dimensions in the $\MS$ scheme. It collects the one-loop self-energy of the emitted gluon, the vertex corrections in which the emitted gluon splits and recombines, and the one-$\rho$ (commutator) parts of the diagrams with two valence vertices (Figures~\ref{fig:vac_g3_onerho} and~\ref{fig:vac_g3_tworho}). The pole multiplies the combination $\frac{11}{3}-2\log(k^+/\Lambda)$, whose first term is $\beta_0/N_c$ for pure glue, $\beta_0=\frac{11}{3}N_c$. In coordinate space,
\begin{align}\label{eq:A3onerhoX}
\Big[\mf A^{(3)a}_{2i}(k^+,-\v z)\Big]_{\rho}=\ &-\int_x\frac{ig^3N_c\,\rho^a(\v x)\,\Kv^i(\v x-\v z)}{16\sqrt2\,\pi^3\sqrt{k^+}}\Bigg\{\Big[-\frac2\epsilon-2\gamma_E-\log\frac{(\v x-\v z)^2\mu^2_{\MS}}{4}\Big]
\notag\\&\hspace{1.5cm}\times\Big[\frac{11}{3}-2\log\frac{k^+}{\Lambda}+\log\Big(1-\frac{k^+}{\vee}\Big)\Big]
\notag\\&\hspace{1.5cm}
-\frac{67}{9}+\frac{\pi^2}{3}+\log^2\Big(1-\frac{k^+}{\vee}\Big)+\log\Big(1-\frac{k^+}{\vee}\Big)\log\frac{\vee}{k^+}\Bigg\}.
\end{align}
The two-$\rho$ part is antisymmetric in colour, $\propto f^{abc}\{\rho^b,\rho^c\}$, and is conveniently written in a mixed representation,
\begin{align}\label{eq:A3tworho}
\Big[\mf A^{(3)a}_{2i}(k^+,-\v z)\Big]_{\rho\rho}=\ &\int d^2p\,d^2k\int_{x,y}e^{i\v k\cdot(\v z-\v y)}e^{i\v p\cdot(\v y-\v x)}\,\frac{ig^3f^{abc}\{\rho^b(\v x),\rho^c(\v y)\}}{32\sqrt{2k^+}\,\pi^5\,\v p^2\,\v k^2}
\notag\\&\times\Bigg[\log\frac{\Lambda}{k^+}\,\frac{\v k^2p^i-(\v k\cdot\v p)\,k^i}{(\v k-\v p)^2}
-2k^i\log\frac{\vee}{\Lambda}
-2k^i\log\frac{k^+\v p^2}{\vee\,\v k^2+k^+\v p^2}
\notag\\&\quad
+\frac{(2\v p^2-2\v k\cdot\v p+\v k^2)\,k^i+\v k^2p^i}{(\v k-\v p)^2}\log\Big[1+\Big(\frac{\vee}{k^+}-1\Big)\frac{\v k^2}{(\v k-\v p)^2}\Big]
\notag\\&
-\frac{\v p^2k^i+\v k^2p^i}{(\v k-\v p)^2}\log\frac{\vee-k^+}{\Lambda}-k^i\log\frac{\vee}{k^+}\Bigg],
\end{align}
where, inside the integral, $\v k$ and $\v p$ are integration variables. The three-$\rho$ part originates from the third-order terms of Eq.~\eqref{eq:lcpt} in which a gluon is emitted and reabsorbed by the valence charges before or after the emission of the observed gluon (Figure~\ref{fig:vac_g3_threerho}). The terms proportional to $1/\Lambda$ and $1/\vee$ cancel exactly, and the remainder is
\begin{align}\label{eq:A3threerho}
\Big[\mf A^{(3)a}_{2i}(k^+,-\v z)\Big]_{\rho\rho\rho}&=-\log\frac{\vee}{\Lambda}\int_{w,x,y,u}\frac{ig^3\,\rho^a(\v w)\,\rho^b(\v x)\,\rho^b(\v y)\,K(\v x,\v y,\v u)\,\Kv^i(\v w-\v z)}{8\sqrt2\,\pi^4\sqrt{k^+}}
\notag\\&=\mf A^{(1)a}_{2i}(k^+,-\v z)\,\mf N_2\,.
\end{align}
The operator ordering in Eq.~\eqref{eq:A3threerho}, with $\rho^a(\v w)$ to the left of the norm factor, is dictated by Eq.~\eqref{eq:lcpt}. It is important in what follows: since $\mf N_2\mf A^{(1)}\neq\mf A^{(1)}\mf N_2$, the opposite ordering would differ by the commutator $[\mf N_2,\mf A^{(1)}]$, which is a two-$\rho$ term of exactly the form and size of the term $-2k^i\log(\vee/\Lambda)$ of Eq.~\eqref{eq:A3tworho} (Eq.~\eqref{eq:commN2A} and the discussion there).

The one-$\rho$, two-$\rho$ and three-$\rho$ parts of the wave function from which $\mf A^{(3)}_2$ is obtained agree with the NLO light-cone wave function of Lublinsky and Mulian~\cite{Lublinsky:2016meo} (eqs.~3.97--3.102 there), including the coefficient $11/3$, the constant $67/9-\pi^2/3$, the dependence on the longitudinal cut-offs and the three-$\rho$ term; the two-$\rho$ parts agree after the antisymmetrisation enforced by $f^{abc}\{\rho^b,\rho^c\}$. The coordinate-space forms~\eqref{eq:A3onerhoX}--\eqref{eq:A3threerho} go beyond the published material.

\paragraph{Two-gluon emission.} The coefficient $\mf B_2$ gives the amplitude for two soft gluons, with longitudinal momenta $k^+$ and $p^+$, at transverse positions $\v w$ and $\v z$:
\begin{align}\label{eq:B2}
\mf B^{bc}_{2ji}(k^+,-\v w;p^+,-\v z)=\ &\int_x\frac{ig^2f^{cba}\,\rho^a(\v x)\,\sqrt{p^+k^+}}{4\pi^2\big[p^+(\v x-\v z)^2+k^+(\v x-\v w)^2\big](p^++k^+)}
\notag\\&\times\Bigg[\frac{\delta^{ij}\big[(\v x-\v z)^2-(\v x-\v w)^2\big]}{2(\v z-\v w)^2}
\notag\\&\quad
+\frac{p^++k^+}{k^+}\Big(\frac{(x-z)^j(z-w)^i}{(\v z-\v w)^2}-\frac{(x-z)^j(x-w)^i}{2(\v x-\v w)^2}\Big)
\notag\\&\quad
+\frac{p^++k^+}{p^+}\Big(\frac{(x-w)^i(z-w)^j}{(\v z-\v w)^2}+\frac{(x-z)^j(x-w)^i}{2(\v x-\v z)^2}\Big)\Bigg]
\notag\\&\quad\times\Theta(\vee-k^+-p^+)
\notag\\&
-\int_{x,y}\frac{g^2\,\{\rho^c(\v y),\rho^b(\v x)\}\,\Kv^i(\v y-\v w)\,\Kv^j(\v x-\v z)}{8\pi^2\sqrt{p^+k^+}}
\notag\\&\quad\times\Theta(\vee-k^+)\,\Theta(\vee-p^+)\,,
\end{align}
with $k^+,p^+>\Lambda$. The first term, linear in $\rho$, arises from the emission of a gluon of momentum $k^++p^+$ which splits into two (triple-gluon vertex), from the instantaneous emission of the pair, and from the commutator part of the emission of two gluons by two charges (Figure~\ref{fig:vac_gg}); the parent gluon is a soft mode, hence $k^++p^+<\vee$. The second term is the independent emission of the two gluons by two valence charges; each gluon is separately a soft mode, and there is no condition on $k^++p^+$.
% CHOICE FOR THE AUTHORS: the support of the commutator part of the two-emission diagram (contained in the first term) is taken as k^+ + p^+ < vee, as for the splitting diagrams; by the kinematics of that diagram it could also be the full square Lambda<k^+,p^+<vee.

\paragraph{Scattering of a soft gluon on the valence charges.} The coefficient $\mf B_3$ describes the transition of a soft gluon of momentum $k^+$ at $\v z$ into a soft gluon of momentum $p^+$ at $\v u$ through its interaction with the valence charges: instantaneous interaction, emission of a second gluon followed by absorption of the first, and absorption of the gluon followed by re-emission (Figure~\ref{fig:g_g}). For $p^+>k^+$ the valence charges transfer longitudinal momentum to the gluon, for $p^+<k^+$ the gluon transfers momentum to the valence charges. With $\Theta_>\equiv\Theta(p^+-k^+-\Lambda)$ and $\Theta_<\equiv\Theta(k^+-\Lambda-p^+)$,
\begin{align}\label{eq:B3}
\mf B^{ba}_{3ji}(k^+,\v z;p^+,-\v u)=\ &\frac{ig^2f^{abc}}{(2\pi)^2}\int_x\frac{\sqrt{p^+k^+}\,(p^++k^+)}{(k^+-p^+)^2}\,\frac{\rho^c(\v x)\,\delta_{ij}}{p^+(\v x-\v u)^2-k^+(\v x-\v z)^2}
\notag\\&\qquad\times\big[\Theta_>+\Theta_<\big]
\notag\\&
+\frac{ig^2f^{abc}}{2\pi^2k^+}\int_x\frac{\sqrt{k^+p^+}}{k^+-p^+}
\notag\\&\qquad\times\frac{\big(k^+(x-z)^m-p^+(x-u)^m\big)\big((x-z)^l-(x-u)^l\big)}{k^+(\v x-\v z)^2-p^+(\v x-\v u)^2}\,\rho^c(\v x)
\notag\\&\qquad\times\Big(\delta_{il}\delta_{jm}-\frac{k^+}{p^+}\delta_{jl}\delta_{im}-\frac{k^+}{k^+-p^+}\delta_{ij}\delta_{lm}\Big)
\notag\\&\qquad\times\Bigg[-\frac{\Theta_>}{(\v u-\v z)^2}-\frac{p^+k^+\,\Theta_<}{\big(k^+(\v x-\v z)-p^+(\v x-\v u)\big)^2}\Bigg]
\notag\\&
-\frac{g^2}{4\pi^2\sqrt{p^+k^+}}\int_{x,y}\Kv^i(\v x-\v z)\,\Kv^j(\v y-\v u)\,\{\rho^a(\v x),\rho^b(\v y)\}
\notag\\&
-\frac{ig^2f^{abc}}{4\pi^2\sqrt{p^+k^+}}\int_x\Kv^i(\v x-\v z)\,\Kv^j(\v x-\v u)
\notag\\&\qquad\times\frac{k^+(\v x-\v z)^2+p^+(\v x-\v u)^2}{p^+(\v x-\v u)^2-k^+(\v x-\v z)^2}\,\rho^c(\v x)\,.
\end{align}
The double pole $1/(k^+-p^+)^2$ at the boundary $|p^+-k^+|=\Lambda$ cancels between the instantaneous and the non-instantaneous terms of the first two lines, so that no term proportional to $k^+/\Lambda$ arises. The vanishing denominators $p^+(\v x-\v u)^2=k^+(\v x-\v z)^2$ are principal values (Section~\ref{sec:fin_conv}).

\paragraph{Gluon splitting.} The coefficient $\mf C_1$ is independent of $\rho$; it is the light-cone amplitude for the splitting of a soft gluon $(k^+,\v z)$ into two soft gluons of momenta $p^+$ and $k^+-p^+$ (Figure~\ref{fig:g_gg}):
\begin{align}\label{eq:C1}
\mf C^{bca}_{1jki}(k^+,\v z;p^+,\v x)=\ &-\frac{gf^{abc}}{2\pi\sqrt{2k^+}}\,\frac{\sqrt{p^+(k^+-p^+)}}{k^+}\,\delta^{(2)}\Big(\v z+\frac{p^+}{k^+}\v x\Big)\,\Kv^m(\v x)
\notag\\&\times\Big[\delta_{im}\delta_{jk}-\frac{k^+}{p^+}\delta_{jm}\delta_{ik}-\frac{k^+}{k^+-p^+}\delta_{km}\delta_{ij}\Big]\Theta(k^+-p^+-\Lambda)\,.
\end{align}
The $\delta$-function expresses the conservation of the transverse centre of mass $p^+\v x_1+(k^+-p^+)\v x_2=k^+\v x_0$ in the splitting.

\paragraph{Coefficients fixed by unitarity.} $\mf A^{(3)}_1$ and $\mf B_1$, which multiply annihilation operators, are not independent; they are given in terms of the coefficients above by Eqs.~\eqref{eq:A31unit} and~\eqref{eq:B1unit}. In the cross section they appear only through these relations.

\begin{summarybox}
\begin{itemize}
\item The soft gluon cloud of the projectile is generated by a unitary operator $\Omega$, Eq.~\eqref{eq:OmegaExp}, whose coefficients are functionals of the non-commuting charge density $\rho$.
\item $\mf N$, $\mf A_2$, $\mf B_2$ follow from the dressed vacuum, $\mf B_3$, $\mf C_1$ from the dressed one-gluon state, the rest from unitarity (Table~\ref{tab:coeffs}).
\item New in this paper: $\mf A^{(3)}_2$ with its one-, two- and three-$\rho$ parts, Eqs.~\eqref{eq:A3onerho}--\eqref{eq:A3threerho}, in coordinate space and with the operator ordering $[\mf A^{(3)}_2]_{\rho\rho\rho}=\mf A^{(1)}_2\mf N_2$.
\item Supports in longitudinal momentum are fixed by kinematics: $k^++p^+<\vee$ where a soft parent gluon occurs, the full square $\Lambda<k^+,p^+<\vee$ for two independent emissions, $|p^+-k^+|>\Lambda$ for momentum exchanged with the valence charges.
\end{itemize}
\end{summarybox}
